\documentclass[UTF8,11pt,fontset=fandol]{ctexart}
\usepackage[a4paper,left=23mm,right=23mm,top=20mm,bottom=21mm]{geometry}
\usepackage{amsmath,amssymb,booktabs,tabularx,array,graphicx,float,enumitem}
\usepackage{xcolor,listings,caption,fancyhdr,hyperref}
\hypersetup{colorlinks=true,linkcolor=black,urlcolor=blue!55!black,pdfauthor={唐骏逸},pdftitle={CacheLab 实验报告}}
\setlength{\parindent}{2em}
\setlength{\parskip}{3pt}
\linespread{1.18}
\setlength{\emergencystretch}{2em}
\setlist[itemize]{leftmargin=2em,itemsep=2pt,topsep=4pt}
\setlist[enumerate]{leftmargin=2em,itemsep=2pt,topsep=4pt}
\captionsetup{font=small,labelfont=bf,skip=5pt}
\ctexset{
  section={format=\Large\bfseries,beforeskip=15pt,afterskip=8pt},
  subsection={format=\large\bfseries,beforeskip=11pt,afterskip=6pt},
  subsubsection={format=\normalsize\bfseries,beforeskip=8pt,afterskip=4pt}
}
\pagestyle{fancy}
\fancyhf{}
\fancyhead[L]{\small CacheLab 实验报告}
\fancyhead[R]{\small 唐骏逸}
\fancyfoot[C]{\small\thepage}
\setlength{\headheight}{14pt}
\renewcommand{\headrulewidth}{0.3pt}
\lstset{
  language=C++,basicstyle=\small\ttfamily,keywordstyle=\bfseries,
  commentstyle=\color{black!60},columns=fullflexible,keepspaces=true,
  showstringspaces=false,breaklines=true,frame=single,
  rulecolor=\color{black!20},xleftmargin=0.5em,xrightmargin=0.5em,
  aboveskip=6pt,belowskip=6pt
}
\newcommand{\mc}{\texttt{miss\_cache}}
\newcommand{\mr}{\texttt{miss\_reg}}
\newcommand{\cmd}[1]{\texttt{\detokenize{#1}}}
\newcommand{\shot}[4][0.96]{%
  \begin{figure}[H]\centering
  \includegraphics[width=#1\linewidth]{#2}
  \caption{#3}\label{#4}
  \end{figure}}
\newenvironment{stageblock}{%
  \par\addvspace{6pt}\noindent\begin{minipage}{\linewidth}
  \setlength{\parindent}{2em}
}{\end{minipage}\par\addvspace{6pt}}
% 所有原始终端截图来自上一轮报告源码包，保存于本报告的 img/ 目录。
% 以下成绩来自已保存的本地评测，本次另用现有 trace 核对 Part B 数值。
\begin{document}
\thispagestyle{empty}
\begin{center}
  {\LARGE\bfseries CacheLab 报告}\\[7pt]
  姓名：唐骏逸\qquad 学号：2025201863\qquad 日期：2026年11月18日
\end{center}
\vspace{-5pt}
\begin{table}[H]
\centering\small
\setlength{\tabcolsep}{4pt}
\begin{tabularx}{\linewidth}{*{6}{>{\centering\arraybackslash}X}}
\toprule
csim 分数 & case1 & case2 & case3 & weighted & Part B 性能分\\
（/100）& speedup & speedup & speedup & speedup &（/100）\\
\midrule
100.00 & 8.24 & 9.28 & 7.06 & 8.08 & 100.00\\
\bottomrule
\end{tabularx}
\end{table}
\vspace{-9pt}
\subsection*{本地评测截图}
{\small 以下为实验过程中保存的原始终端截图，展示 Part A 完整测试与 Part B 最终测试结果。}

\vspace{-7pt}
\shot[0.98]{img/part_a.png}{Part A 本地评测：20/20，通过分数100.00。}{fig:parta}
\vspace{-7pt}
\shot[0.98]{img/part_b_final.png}{最终 Part B 本地评测：综合加速比8.0797，性能分100.00。}{fig:final}

\clearpage


\section{Part A: cache 模拟器}
\subsection{实现简述}
这一部分我先写了每组只有一行的直接映射版本，确认地址拆分、命中和替换的计数没有问题后，再扩展到每组 \(E\) 行，并加入 LRU。缓存用一个连续数组保存：组数为 \(S=2^s\)，块大小为 \(B=2^b\)，第 \(q\) 组从 \cmd{cache[q*E]} 开始。每行记录 \texttt{valid}、\texttt{tag} 和上次访问的序号 \texttt{last\_used}。

\textbf{地址拆分。} 对地址 \(a\)，先右移 \(b\) 位去掉块内偏移，再取组号和 tag：
\[
  u=a\gg b,\qquad
  q=u\mathbin{\&}(2^s-1),\qquad
  \mathrm{tag}=u\gg s.
\]
判断命中时，要先看 \texttt{valid}。初始化后的空行即使 tag 恰好相同，也不能算作命中。

\textbf{查找与替换。} 每次访问时，访问序号加1，再扫描目标组。如果找到相同 tag 的有效行，就增加 \texttt{hits}，并更新它的访问序号。没有命中则增加 \texttt{misses}，先找空行装入；没有空行时，才替换 \texttt{last\_used} 最小的行，并增加 \texttt{evictions}。装入新块后也要更新访问序号。扫描中遇到空行不能直接停下，因为后面的有效行仍有可能命中。

L 和 S 都按这个过程处理，写缺失时也会装入缓存块。这里没有保存块中的实际数据，也没有统计脏块的回写流量，所以代码中未设置 dirty 位；实现的是本题要求的三项计数。

\textbf{手工核对。} 调试 LRU 时，我用较短的 \cmd{yi.trace}，取 \(s=4,E=2,b=4\)，逐条检查组号、tag 和替换对象。下表中的时间就是访问序号，只列出本次访问涉及的组。
\begin{table}[H]
\centering\small
\setlength{\tabcolsep}{5pt}
\begin{tabular}{clccp{64mm}}
\toprule
次序 & 访问 & 组号 & tag & 结果与访问后的目标组\\
\midrule
1 & 读 \texttt{0x10}  & 1 & 0 & miss；\(\{0@1,\text{空}\}\)\\
2 & 读 \texttt{0x20}  & 2 & 0 & miss；\(\{0@2,\text{空}\}\)\\
3 & 读 \texttt{0x22}  & 2 & 0 & hit；\(\{0@3,\text{空}\}\)\\
4 & 写 \texttt{0x18}  & 1 & 0 & hit；\(\{0@4,\text{空}\}\)\\
5 & 读 \texttt{0x110} & 1 & 1 & miss；\(\{0@4,1@5\}\)\\
6 & 读 \texttt{0x210} & 1 & 2 & miss + eviction；\(\{2@6,1@5\}\)\\
7 & 读 \texttt{0x12}  & 1 & 0 & miss + eviction；\(\{2@6,0@7\}\)\\
\bottomrule
\end{tabular}
\end{table}
结果是 \texttt{hits:2 misses:5 evictions:2}。第6次访问淘汰了 tag 0，所以第7次再访问它时仍然缺失。另外，一次替换也要计入一次缺失，两个计数不能二选一。完成这些检查后，运行 \cmd{python3 test/csim_test.py}，20组测试均与参考程序一致，结果见前面的截图。

\subsection{亮点}
\begin{itemize}
  \item \textbf{把映射和替换分开调试。} 先检查地址拆分和 \(E=1\) 的情况，再加入组内扫描与 LRU。出现错误时，可以先判断是找错了组，还是选错了替换行。
  \item \textbf{用访问序号记录 LRU。} 命中和装入时都更新序号，替换时找最小值即可。每次只扫描目标组的 \(E\) 行，时间复杂度为 \(O(E)\)，缓存数组占用 \(O(2^sE)\) 的空间。
\end{itemize}

\section{Part B: 矩阵乘法优化}
\subsection{亮点}
\begin{enumerate}
  \item \textbf{寄存器分块：}case1、case2 一次计算 \(4*4\) 个结果，让读入的 A、B 元素多用几次，部分和留在寄存器中，算完一块再写回 C。
  \item \textbf{为 case3 调整块大小：}改用 \(3*7\) 分块，27行恰好分成9组，最后不足7列的部分单独判断边界。
  \item \textbf{调整 case3 的遍历顺序：}先固定列块，再遍历行块，相邻块交替正向、反向扫描 \(k\)，尽量接着使用刚访问过的 B。
  \item \textbf{同时比较两项访存开销：}除了 cache miss，还统计 load/store 次数，最后按 latency 判断修改是否有效。
\end{enumerate}

\subsubsection{我认为的最优秀的实现排序}
\begin{enumerate}
  \item \textbf{case3：}没有继续沿用统一的 \(4*4\) 方案，而是结合27行、27列的尺寸调整块形状和遍历顺序，解决了最后一个明显拖慢综合成绩的 case。
  \item \textbf{case1：}从交换循环到整行累加，再到 \(4*4\) 分块，每一步都有结果可对照，也最能看出只减少 cache miss 还不够。
  \item \textbf{case2：}直接沿用了 case1 的分块方法，代码改动不多，但加速比达到了9.2808。
\end{enumerate}

\paragraph{实验设定与指标。}
A、B、C 都按行存储，形状分别为 \(m*n\)、\(n*p\)、\(m*p\)，每个元素占4字节。模拟缓存共有32组，每组1行，每块16字节，总容量512字节；A、B、C 在内存中紧接着排列。代码中的局部变量使用框架提供的 \texttt{reg}，最多可同时使用36个模拟寄存器，没有另申请内存，也没有使用 buffer。下文说的寄存器数组，都是指框架里的 \texttt{reg} 数组。

评测中的 \(M=\mc\) 表示缓存缺失次数，\(R=\mr\) 统计的是 load/store 次数。后一个名字容易误解，它并不是另一级缓存的缺失次数。延迟和加速比按下面的公式计算：
\[
  L=15M+R,\qquad S_i=\frac{L_{i,\mathrm{baseline}}}{L_i},\qquad
  S_{\mathrm{weighted}}=0.3S_1+0.3S_2+0.4S_3.
\]
测试修改后的代码用 \cmd{python3 test/gemm_test.py}，带 \cmd{--baseline} 时只测原始基准。下面比较的版本都通过了正确性检查。

\subsection{case1}
\begin{stageblock}
\textbf{形状与起点。} A、B、C 均为 \(16*16\)。最初使用 \texttt{ijk} 循环，内层沿 B 的一列读取，相邻地址相差64字节。虽然一个缓存块能放4个元素，但这种读法很难连续用完它们。原始代码把单个输出的部分和存在寄存器里，C 只写回一次，因此我先考虑改善 B 的读取顺序。
\shot[0.94]{img/stage_4.png}{原始 \texttt{ijk} 版本。三个 case 的 baseline latency 分别为80328、607040、613542。}{fig:baseline}
\end{stageblock}

\begin{stageblock}
\textbf{第一步：交换循环。} 我先把循环改为 \texttt{ikj}，让 B、C 都按行访问。cache miss 确实从4792降到了2464，但加速比只有1.6143。继续检查访存次数，才发现每次乘加都要读写 C，load/store 反而从8448增到了12800。这一步改善了 B 的局部性，却增加了 C 的访问开销。测试时还误用过 \cmd{--baseline}，当时修改后的代码根本没有被测到，改正命令后才得到下面的结果。
\shot[0.91]{img/stage_7.png}{改为 \texttt{ikj} 后，cache miss 减少，但 C 的反复读写增加了访存次数。}{fig:ikj}
\end{stageblock}

\begin{stageblock}
\textbf{第二步：整行寄存器累加。} 既然多出的访存来自 C，我就用16个寄存器保存一整行的部分和，等全部 \(k\) 算完再写回。这样，B 仍然可以按行读取，中途也不用反复访问 C。总访存次数为
\[
  R=16^2+16^3+16^2=4608.
\]
这次加速比提高到了3.6106。不过，换到下一行 C 时，B 还要重新读取，所以接下来需要考虑多行之间的复用。中间也讨论过2列、4列展开，但没有单独保留这些方案的测试记录，因此表中只列实际测过的版本。
\shot[0.91]{img/stage_8.png}{一整行的部分和留在寄存器中，算完再写回，加速比提高到3.6106。}{fig:row}
\end{stageblock}

\begin{stageblock}
\textbf{第三步：同时复用 A、B。} 接着改为一次计算 \(4*4\) 个输出，用16个寄存器保存部分和。每轮 \(k\) 读取4个 A、4个 B，完成16次乘加。这样，A 的每个元素能用于4列，B 的每个元素能用于4行，最后再把这一块 C 写回。16个输出块的总访存次数为
\[
  R=16\,[16(4+4)+16]=2304,
\]
与实测相同。此时 cache miss 为496，latency 为9744，加速比达到了8.2438。
\shot[0.91]{img/stage_9.png}{此时只修改了 case1，单项加速比为8.2438，综合加速比为3.1732。}{fig:tile}
\end{stageblock}

\begin{table}[H]
\centering\small
\begin{tabular}{lrrrr}
\toprule
case1 实测版本 & \mc & \mr & latency & speedup\\
\midrule
原始 \texttt{ijk} & 4792 & 8448 & 80328 & 1.0000\\
交换为 \texttt{ikj} & 2464 & 12800 & 49760 & 1.6143\\
整行寄存器累加 & 1176 & 4608 & 22248 & 3.6106\\
\(4*4\) 寄存器分块 & 496 & 2304 & 9744 & 8.2438\\
\bottomrule
\end{tabular}
\caption{case1 各版本的本地评测结果。}
\end{table}

\textbf{cache miss 的计算。} 访存次数可以直接数出来，但 cache miss 还要看这些访问落在哪一组。三个矩阵各占64个缓存块，第一次读写这些块至少需要192次缺失。实测为496次，多出的304次来自块的重复装入，不能简单把2304次访存除以4来估计。按数组的实际基址，A、B 的组号分别为：
\[
 q_A(i,k)=\bigl(4i+\lfloor k/4\rfloor\bigr)\bmod32,\qquad
 q_B(k,j)=\bigl(4k+\lfloor j/4\rfloor\bigr)\bmod32.
\]
例如，A 的第0行第0列和 B 的第0行第0列都落在第0组，但 tag 不同，会互相替换。即使一个块里的元素还没用完，也可能已经被另一个矩阵的数据挤出去。

按照源码的访问顺序逐次生成地址，并记录32组中各自保存的 tag，可以得到下表。每一格对应一个输出块，包括读取 A、B 和写回 C 引起的缺失。换到下一个输出块时，保留上一块计算结束后的缓存状态。
\begin{table}[H]
\centering\small
\begin{tabular}{crrrr}
\toprule
行块起点 $i$ / 列块起点 $j$ & 0 & 4 & 8 & 12\\
\midrule
0  & 39 & 28 & 31 & 28\\
4  & 36 & 27 & 28 & 31\\
8  & 39 & 28 & 31 & 28\\
12 & 36 & 27 & 28 & 31\\
\bottomrule
\end{tabular}
\caption{按地址顺序计算的 case1 各输出块 cache miss。}
\end{table}
合计496次，其中 A、B、C 分别占176、256、64次，与保存的 trace 一致。这张表是按完整地址序列回放得到的，也包含了不同输出块之间的缓存影响。复核代码放在 \cmd{report/verify_trace.py} 中，在项目根目录运行 \cmd{python3 report/verify_trace.py} 即可。脚本还会核对现有8份 trace 的地址顺序和访存数据；它用于复查记录，不能代替官方脚本的正确性检查。

\subsection{case2}
\textbf{形状与实现。} A、B、C 均为 \(32*32\)。只改完 case1 时，综合加速比还只有3.1732，所以我把已经测过的 \(4*4\) 分块方法也用到了 case2。32能被4整除，输出恰好分为64个块，不需要另外处理尾块。

\textbf{理论访问数。} 每块计算32轮 \(k\)，每轮读8个元素，最后写回16个结果，因此
\[
  R=64\,[32(4+4)+16]=17408.
\]
实测 \(\mc=3200\)、\(\mr=17408\)，\(L=65408\)。原始版本的 cache miss 为36032，latency 为607040，最终加速比为9.2808。

\textbf{cache miss。} B 的一行占128字节，原先按列访问时，相邻两次读取相隔较远。改成连续读4列后，一个缓存块里的4个元素可以接着用完。不过，三个矩阵共有768个不同的块，缓存只能放32块，计算后面的输出块时仍然要重新装入不少数据，因此实测3200次 miss 并不接近冷缺失下界。case2 沿用同一套方法已经有较好的效果，后面主要继续修改 case3。

\subsection{case3}
\begin{stageblock}
\textbf{形状与瓶颈。} A 为 \(27*41\)，B 为 \(41*27\)，C 为 \(27*27\)。三个 case 都换成 \(4*4\) 分块后，case3 的加速比只有5.7176，低于前两个 case。检查尺寸时可以发现，27不能被4整除，B 每行的108字节也不是16字节块大小的整数倍，连续读4列有时会跨过两个缓存块。这个 case 又占综合指标40\%的权重，所以接下来重点改它。
\shot[0.96]{img/stage_10.png}{三个 case 都使用 \(4*4\) 分块时，综合加速比为7.5445，性能分为93.49。}{fig:alltile}
\end{stageblock}

\textbf{块形状。} case3 最后改用了 \(3*7\) 分块，用21个寄存器保存部分和。27行正好分成9组，27列分为 \(7+7+7+6\)，只需判断最后一组的列边界。每轮读取3个 A、最多7个 B，再更新对应的输出。对于 \(a*b\) 的输出块，不额外暂存数据时，总访存次数为
\[
  R=mn\left\lceil\frac{p}{b}\right\rceil
    +np\left\lceil\frac{m}{a}\right\rceil+mp.
\]
前两项分别是 A、B 的读取次数，最后一项是 C 的写回次数。代入两种块尺寸：
\[
 \begin{aligned}
 R_{4*4}&=27\cdot41\cdot7+41\cdot27\cdot7+729=16227,\\
 R_{3*7}&=27\cdot41\cdot4+41\cdot27\cdot9+729=15120.
 \end{aligned}
\]
换成 \(3*7\) 后，B 的读取次数从7749增到了9963，但 A 从7749降到了4428，总计少了1107次访存。这里7列的作用主要是减少 A 的重复读取，并不意味着它与缓存块对齐；cache miss 仍需按实际访问顺序检查。

\textbf{块间复用。} 这一版还调整了遍历顺序：外层先遍历列块 \(j\)，内层遍历行块 \(i\)，使相邻输出块读取 B 的同一段列。相邻块的 \(k\) 方向交替变化，前一块从0到40，下一块就从40到0。这样，下一块开头访问的 B 区间与上一块收尾时相同，尚未被替换的数据便有机会接着使用。方向由 \((i/3+j/7)\bmod2\) 决定，下面只保留这部分循环：
\begin{lstlisting}
for(reg j=0;j<p;j+=7){
  for(reg i=0;i<m;i+=3){
    // 先把当前块的21个部分和清零
    for(reg k=((i/3+j/7)%2?n-1:0);
        k>=0 && k<n;
        k+=((i/3+j/7)%2?-1:1)){
      // 读入A的3个元素和B当前的至多7列
      // 将乘积加到对应的部分和中
    }
    // 当前块算完后，再把结果写回C
  }
}
\end{lstlisting}
虽然 \(k\) 的方向变了，但每个输出仍然累加全部41项，没有漏算或重复计算。最终各 case 的正确性检查均为 \texttt{OK}。

\begin{table}[H]
\centering\small
\begin{tabular}{lrrrr}
\toprule
case3 实测版本 & \mc & \mr & latency & speedup\\
\midrule
朴素 \texttt{ijk} & 36869 & 60507 & 613542 & 1.0000\\
统一 \(4*4\) 分块 & 6072 & 16227 & 107307 & 5.7176\\
\(3*7\) + 块次序 + 双向 \(k\) & 4789 & 15120 & 86955 & 7.0559\\
\bottomrule
\end{tabular}
\caption{case3 的测试结果。最后一版同时改了块形状、块的遍历顺序和 \(k\) 的方向。}
\end{table}

\textbf{理论与实测。} 最终的15120次 load/store 与上面的计算相同，cache miss 为4789次。相比 \(4*4\) 版本，miss 少了1283次，latency 下降约18.97\%。此前按地址顺序计算得到的 miss 也是4789，之后用官方脚本测试得到了相同结果；再次回放现有的 \cmd{gemm_traces/case3.trace}，数值也一致。不过，这一版把三处修改放在了一起，没有分别做测试，因此无法从这组结果中确定每项修改单独带来了多少收益。

\textbf{综合结果。} case1、case2 保持原来的实现，只替换 case3 后，综合加速比为
\[
 S_{\mathrm{weighted}}
 =0.3\frac{80328}{9744}
  +0.3\frac{607040}{65408}
  +0.4\frac{613542}{86955}
 \approx 8.0797.
\]
综合性能分也从93.49提高到了100.00，最后保留了这一版。三个 case 的尺寸不同，分别选择合适的分块方式，比一直使用同一个块尺寸更有效。

\section{反馈/收获/感悟/总结}
这次实验中，交换循环的结果最值得回看。cache miss 已经明显减少，加速比却没有提高很多，原因是 C 的读写次数增加了。如果只看 miss，就会忽略这部分开销。后面改用寄存器保存部分和，再同时计算多行、多列，才逐步减少了总访存次数。这几步的测试结果放在一起，能比较清楚地看出每次修改究竟解决了什么问题。

case3 则让我注意到，已经在前两个尺寸上测过的方法，换一个尺寸还需要重新检查。\(4*4\) 用在16列、32列时很方便，到了27列就会留下尾块，B 的行跨度也变了。选择 \(3*7\) 时，需要同时算 A、B 各自会读多少次，还要留够寄存器给循环变量和临时值。按当前代码中变量的同时存活情况，\(4*4\) 最多使用30个模拟寄存器，\(3*7\) 最多33个，都在36个的限制内；这个数是按源码统计的，运行时仍由框架检查。

测试命令也不能忽略。误用 \cmd{--baseline} 时，改了代码却一直在测原始版本，当然看不到变化。此后比较版本时，除了检查结果是否正确，也会先确认命令实际测的是哪一份实现。保留中间版本的截图很有用：只留下最终的满分结果，很难再说明交换循环、整行累加和分块各自带来了什么变化。

\section{参考的重要资料}
\begin{enumerate}
  \item Randal E. Bryant, David R. O'Hallaron. \textit{Computer Systems: A Programmer's Perspective}, 3rd ed. 第6章的存储层次、局部性、缓存映射及矩阵计算示例。
\end{enumerate}

\end{document}

